\documentclass[11pt]{article}
\usepackage[T1]{fontenc}
\usepackage{textcomp}
\usepackage[margin=0.75in]{geometry}
\usepackage{amsmath,amssymb,amsthm}
\usepackage{array,booktabs}
\usepackage{hyperref}
\setlength{\parindent}{0pt}
\newtheorem{theorem}{Theorem}
\theoremstyle{definition}
\newtheorem{definition}{Definition}
\title{Flow Proof Artifact: prop-07-triangles-on-same-base}
\author{Generated by \texttt{flow doc proof}}
\date{Flow Proof Book}
\begin{document}
\maketitle
On the same base and on the same side, equal-sided triangles cannot have distinct vertices.\\[0.5em]
\textbf{Source.} euclid — Elements, Book I, Proposition 7\\[0.5em]
\bigskip
\noindent{\large \textbf{Derived fact 1.} \textit{On the same base and on the same side, equal-sided triangles cannot have distinct vertices} \hfill the Euclidean plane \cdot Euclid Book I \cdot Proposition 7: two triangles on the same base cannot meet above it}\par
\smallskip\noindent\textit{Coordinate: the Euclidean plane \cdot Euclid Book I \cdot Proposition 7: two triangles on the same base cannot meet above it}\par
\smallskip\noindent\textit{Source: euclid — Elements, Book I, Proposition 7}\par
\smallskip\noindent\textit{Built on: proposition 4: side-angle-side congruence, for Book I of the Elements in on the Euclidean plane, proposition 5: base angles of an isosceles triangle are equal, for Book I of the Elements in on the Euclidean plane}\par
\medskip
\noindent\textbf{Goal.} On the same base and on the same side, equal-sided triangles cannot have distinct vertices\par
\noindent$\displaystyle \text{triangles on base AB unique}$\par
\medskip
\noindent
\begin{tabular}{@{} >{\raggedright\arraybackslash}p{0.06\textwidth} >{\raggedright\arraybackslash}p{0.47\textwidth} >{\raggedleft\arraybackslash}p{0.41\textwidth} @{}}
\textbf{\#} & \textbf{Proof} & \textbf{Mathematics} \\
\midrule
\textbf{1.} & We prove that proposition 7: two triangles on the same base cannot meet above it for Book I of the Elements in the Euclidean plane. & \\[0.45em]
\textbf{2.} & We invoke the derived fact: segment\_AB == segment\_AC. & \\[0.45em]
\textbf{3.} & We invoke the derived fact: segment\_DB == segment\_DC. & \\[0.45em]
\textbf{4.} & We invoke the derived fact: points\_C\_and\_D\_lie\_on\_same\_side\_of\_line\_AB. & \\[0.45em]
\textbf{5.} & We invoke the axiom governing Book I of the Elements in the Euclidean plane: proposition 4: side-angle-side congruence, for Book I of the Elements in on the Euclidean plane. & \\[0.45em]
\textbf{6.} & We invoke the derived fact governing Book I of the Elements in the Euclidean plane: proposition 5: base angles of an isosceles triangle are equal, for Book I of the Elements in on the Euclidean plane. & \\[0.45em]
\textbf{7.} & From step 2, step 3, step 4, step 5, and step 6, we can deduce that point C equals point D. & $\displaystyle \text{point C} = \text{point D}$ \\[0.45em]
\textbf{8.} & We can deduce that triangles on base AB unique. Hence proven. & $\displaystyle \text{triangles on base AB unique}$ \\[0.45em]
\bottomrule
\end{tabular}
\label{thm:Geometry-Euclid Book I-proposition_7_two_triangles_on_the_same_base_cannot_meet_above_it}
\par\medskip\hrule
\end{document}
